Dựa trên kết quả mô phỏng bằng phương pháp RK4, thấy rõ sự khác biệt giữa chuyển động của đạn pháo trong trường hợp có và không có lực cản không khí.

Cụ thể, trong trường hợp không có lực cản, quỹ đạo đạn có dạng parabol lý tưởng, tầm bắn và độ cao cực đại lớn hơn đáng kể. Góc phóng khoảng \qty{45}{\degree} cho tầm xa nhất, phù hợp với lý thuyết cổ điển. Ngược lại, khi xét đến lực cản không khí với các hệ số cản khác nhau, quỹ đạo bị nén xuống, độ cao và tầm bắn đều giảm mạnh. Hệ số cản càng lớn thì đạn càng rơi sớm và quỹ đạo càng dốc.

Ngoài ra, kết quả cũng cho thấy góc phóng tối ưu trong trường hợp có lực cản không còn chính xác là \qty{45}{\degree}, mà bị lệch đi tùy theo độ lớn của lực cản. Điều này phản ánh tính phi tuyến và phức tạp của bài toán thực tế.

Tóm lại, phương pháp RK4 đã chứng minh hiệu quả trong việc mô phỏng chính xác chuyển động đạn pháo, đồng thời giúp làm rõ sự ảnh hưởng của lực cản không khí đến quỹ đạo chuyển động đạn pháo. Đây là cơ sở để mở rộng nghiên cứu với các mô hình thực tế hơn trong tương lai.
